\paragraph{Uniform comparison on polynomial ranges of dyadic multipliers.}
The positive multiset (P-F56u) gives a rigorous part of the
remaining estimate (P-F56n). This does not bound the entire
remainder. Let $n\ge5$, $l=n-1$, $M=2^l$, and let $J$ be an
odd integer with $1\le J<M$. Put
\[
 a_J\equiv J(3g)^{-1}\pmod M,\qquad
 Z_{n,J}=A_n(\zeta^{a_J}),\qquad
 R_J=\lceil\log_2(16J)\rceil.
 \tag{P-F56ak}
\]
For negative subscripts use $a_{-J}\equiv-a_J\pmod M$;
the rational Hilbert series gives
$Z_{n,-J}=\overline{Z_{n,J}}$.
Then the following explicit uniform error bound holds:
\[
 \left|\log_2\frac{|Z_{n,J}|}{|Z_{n,1}|}
       +\binom{n-1}{3}\log_2J
       -\mathbf1_{3\mid J}\binom n3\right|
       \le 100n^2R_J^2.
 \tag{P-F56al}
\]
In particular, for every fixed real $K\ge0$, uniformly for
odd $3\le J\le n^K$,
\[
 \log_2\frac{|Z_{n,J}|}{|Z_{n,1}|}
 \le -\frac{\log_2 3-1}{6}n^3
                         +O_K(n^2\log^2 n).
 \tag{P-F56am}
\]
Consequently the sum of the magnitudes of all terms indexed
by $a\equiv\pm J(3g)^{-1}$ with $3\le J\le n^K$ is
$o(|Z_{n,1}|)$. Thus every fixed polynomial range of
multipliers is negligible in (P-F56n). This leaves the
remaining, unrestricted multipliers uncontrolled.

\emph{Sparse dyadic separation.}
We first give the elementary estimate that prevents a high
dyadic block from being cancelled by a very small tail.
Let $1\le r\le3$, $1\le q_1<\cdots<q_r$ be integers,
$\epsilon_i\in\{1,-1\}$, and $\kappa\in\mathbb Z$.
For odd $J\ge1$ and $R=\lceil\log_2(16J)\rceil$, set
\[
 x=\frac\kappa3+
              \frac J3\sum_{i=1}^r\epsilon_i2^{-q_i}.
\]
Then
\[
 \|x\|_{\mathbb R/\mathbb Z}
       \ge\frac1{6\,2^{q_1+(r-1)R}}.
 \tag{P-F56an}
\]
To prove this, retain the initial block ending just before
the first gap $q_{i+1}-q_i>R$, or retain all terms if no
such gap occurs. If the retained block ends at $q_t$, its
numerator over the denominator $3\cdot2^{q_t}$ is odd:
the final term is odd, all earlier terms are even, and
$\kappa2^{q_t}$ is even. It is therefore not an integer
and its distance from an integer is at least
$1/(3\cdot2^{q_t})$.
If a tail remains, its absolute value is less than
\[
 \frac{2J}{3}\,2^{-q_{t+1}}
 \le\frac{J2^{-R}}{3\,2^{q_t}}
 \le\frac1{48\,2^{q_t}}.
\]
The reverse triangle inequality for distance to the
integers now gives at least $1/(6\cdot2^{q_t})$.
Since every retained gap is at most $R$, one has
$q_t\le q_1+(r-1)R$, proving (P-F56an).
In particular, if $q_1\le R$ and $r\le3$, then
\[
 2^{-3R-1}\le |1-e^{2\pi ix}|\le2.
 \tag{P-F56ao}
\]
Here we used $|1-e^{2\pi ix}|\ge4\|x\|$.

\emph{Partition of the exact product.}
Write $w_i=(-2)^i$ for $0\le i<l$, with the remaining
weight zero. Let $b_J$ be the residue $J3^{-1}$ modulo
$M$. Then $\eta_J=\zeta^{b_J}$ satisfies
$\eta_J^3=\zeta^J$, and
\[
 \eta_J=\omega_J\exp\left(\frac{2\pi iJ}{3M}\right),
 \qquad \omega_J^3=1,
 \qquad \omega_J=1\ \Longleftrightarrow\ 3\mid J.
 \tag{P-F56ap}
\]
Indeed $3b_J=J+k_JM$ with $k_J\in\{0,1,2\}$, and
$\omega_J=e^{2\pi i k_J/3}$; because $J<M$, the only
case $k_J=0$ is $3\mid J$.
Every $w_i$ is congruent to one modulo three. The exact
denominator in (P-F56v) therefore splits into:
\begin{enumerate}
\item the $l$ singleton factors $1-\eta_J^{w_i}$;
\item the distinct active triple factors
      $1-\eta_J^{w_i+w_j+w_k}$, $i<j<k<l$,
      omitting $\{0,1,k\}$ and $\{0,2,3\}$;
\item the factors $1-\eta_J^{w_i+w_k}$,
      $0\le i<k<l$, omitting the pair $\{0,1\}$.
\end{enumerate}
The last family consists precisely of triples using the
zero-weight variable. The common factor $1/M$ in the
normalized residues cancels from their ratio.

For any factor let $k$ be the largest active index and put
$q=l-k$. Call it low if $q>R$ and high if $q\le R$,
where $R=R_J$. A sum of at most three distinct active
weights divided by $M$ has the form
$\sum\epsilon_i2^{-q_i}$ with smallest denominator
exponent $q_1=q$. Its absolute value lies between
$2^{-q}/4$ and $2\cdot2^{-q}$ when it has three terms;
for two terms the sharper bounds are
$2^{-q}/2$ and $3\cdot2^{-q}/2$.
All phases are consequently covered by (P-F56an).

\emph{Active triples.}
For an active triple, its integer weight is divisible
by three, so its phase at $J$ is $J$ times its phase
at $1$, with no cube-root shift. On the low factors
the absolute phase is at most
$(2J/3)2^{-q}<1/24$.
The elementary inequalities
$1-y^2/6\le\sin y/y\le1$ show, for $|x|\le1/24$,
\[
 \left|\log_2\frac{\sin(\pi x)}{\pi x}\right|
                                      \le8x^2.
\]
Thus the logarithm of the ratio of the $J$ factor to
the $1$ factor differs from $\log_2 J$ by at most
$16J^2 4^{-q}$. For each $q$ there are at most
$n^2/2$ triples. Summing the geometric series over
$q>R$ gives total error less than $n^2$, since
$2^R\ge16J$.

There are at most $Rn^2/2$ high triples.
Their $J$ factors satisfy (P-F56ao). Their reference
factors have absolute value at least $2^{-q}/3$:
the reference phase has absolute value between
$2^{-q}/12$ and $2\cdot2^{-q}/3$, is less than $1/2$,
and $|1-e^{2\pi ix}|\ge4|x|$ on this interval.
The absolute logarithm of each factor ratio is therefore
at most $5R$. Replacing the number of low triples by
$\binom l3$ costs at most another $Rn^2\log_2 J/2$.
The omitted active triples number $l-1$; each ratio
has absolute logarithm at most $5R$ (or at most $2R$
in the low case). Combining these bounds, their net
contribution to the logarithm of the residue ratio is
\[
               -\binom l3\log_2 J+E_3,
                         \qquad |E_3|\le10n^2R^2.
 \tag{P-F56aq}
\]
The loose constant also covers $R\ge l$, when there
are no low factors.

\emph{Pairs using the zero-weight variable.}
At the reference embedding the phase of a pair is
$1/3$ or $2/3$, plus a perturbation of absolute value
at most $1/4$. Hence every reference factor has
absolute value between $1/3$ and $2$.
If $3\nmid J$, the low target factors also have a
nontrivial cube-root shift, perturbed by less than
$1/24$; they and the reference factors have logarithms
bounded by absolute constants.
There are at most $nR$ high pairs. Applying the
two-term version of (P-F56an) bounds the absolute
logarithm of each high ratio by $3R$.
Their total contribution is therefore bounded by
$4n^2R^2$ in absolute value.

If $3\mid J$, the low target phase has no cube-root
shift. Its factor lies between
$(2/3)J2^{-q}$ and $4J2^{-q}$.
The logarithm of the reciprocal-factor ratio is
therefore $q-\log_2 J$ with error at most four.
For the full pair family,
\[
 \sum_{0\le i<k<l}(l-k)
      =\sum_{k=1}^{l-1}k(l-k)
      =\binom{l+1}{3}=\binom n3.
\]
Discarding the $-\log_2 J$ terms costs at most
$n^2\log_2 J/2$. The high pairs number at most $nR$,
have $q\le R$, and their logarithmic ratios are at
most $3R$ in absolute value. Omitting the pair
$\{0,1\}$ costs at most $n+\log_2 J+4$ in the low
case, or $3R$ in the high case. All these bounds give
the contribution
\[
           \mathbf1_{3\mid J}\binom n3+E_2,
                         \qquad |E_2|\le8n^2R^2.
 \tag{P-F56ar}
\]

\emph{Singletons and conclusion.}
Reference singleton factors have absolute value at
least $2/3$. If $3\nmid J$, the low target factors
stay a fixed distance from one; the at most $R$
high factors satisfy the one-term separation estimate.
If $3\mid J$, each low logarithmic ratio has absolute
value at most $q+\log_2 J+3$, and the high ratios have
absolute value at most $R+2$.
Summing over the $l$ singletons in either case gives
absolute contribution at most $5n^2R^2$.
Together with (P-F56aq)--(P-F56ar), this proves
(P-F56al), with room in the stated constant $100$.

For $J\le n^K$, $R=O_K(1+\log n)$.
Since $J\ge3$, the main terms in (P-F56al) are at most
$-(\log_2 3-1)n^3/6+O_K(n^2\log n)$.
This proves (P-F56am). Complex conjugation gives equal
magnitudes for the multipliers $J$ and $-J$.
For all sufficiently large $n$ the range $J\le n^K$
is smaller than $M/2$, so these indices are distinct.
There are at most $2n^K$ terms, whence
\[
 \limsup_{n\to\infty}\frac1{n^3}\log_2
 \left(
 \frac{\displaystyle\sum_{\substack{3\le J\le n^K\\J\text{ odd}}}
       (|Z_{n,J}|+|Z_{n,-J}|)}{|Z_{n,1}|}
 \right)
 \le-\frac{\log_2 3-1}{6}<0.
 \tag{P-F56as}
\]
An empty sum is interpreted as zero, with logarithm
$-\infty$. This proves the claimed negligible part
of the actual remainder. No comparison of the other
multipliers, and hence no unconditional strict-index
lower bound, follows from this argument.
\hfill$\square$

\paragraph{An exponentially long range with a quantitative bound.}
The explicit error in (P-F56al) gives more than polynomial
ranges. For every odd integer $3\le J\le2^{n/100000}$,
\[
       \log_2\frac{|Z_{n,J}|}{|Z_{n,1}|}
                      \le-\frac{n^3\log_2 J}{36}.
 \tag{P-F56at}
\]
The range may be empty in small dimensions. In its stated
range it satisfies $J<M$ for every $n\ge5$.
For a proof put $t=\log_2J>3/2$. The elementary bounds
$\binom{n-1}{3}\ge n^3/6-n^2$ and
$\binom n3\le n^3/6$, followed by (P-F56al), give
\[
 \log_2\frac{|Z_{n,J}|}{|Z_{n,1}|}
 \le-\frac{(t-1)n^3}{6}
                      +n^2\bigl(t+100(t+5)^2\bigr).
\]
Here $R_J\le t+5$. Since $t-1\ge t/3$ and
$t+5\le5t$, the last bound is at most
$-tn^3/18+2501n^2t^2$.
Now $t\le n/100000$ and $2501/100000<1/36$ prove
(P-F56at). Summing and using conjugation yields the
fully explicit bound
\[
 \frac{\displaystyle\sum_{\substack{3\le J\le2^{n/100000}\\
                                  J\text{ odd}}}
       (|Z_{n,J}|+|Z_{n,-J}|)}{|Z_{n,1}|}
 \le 2^{\,1+n/100000-(\log_2 3)n^3/36}.
 \tag{P-F56au}
\]

There is also a sharper asymptotic range. For every fixed
$0<\delta<1/600$, the conclusion (P-F56as), with the
same negative upper bound, holds when $n^K$ is replaced
by $2^{\delta n}$. We spell out the uniformity argument.
Consider any sequence of admissible odd $J=J_n$ and put
$t_n=\log_2J_n$. If $(t_n)$ is bounded, pass to a
subsequence with the integer $J_n$ constant. Dividing
(P-F56al) by $n^3$ gives upper limit
$(-\log_2J+\mathbf1_{3\mid J})/6$, which is at most
$-(\log_2 3-1)/6$ for $J\ge3$.
If $(t_n)$ is unbounded, pass to a subsequence on which
$t_n\to\infty$. The displayed estimate above, divided
by $n^3$, is
\[
 -\frac{t_n-1}{6}
        +\frac{t_n+100(t_n+5)^2}{n}
 \le-\bigl(1/6-100\delta-o(1)\bigr)t_n+O(1),
\]
because $t_n\le\delta n$. This tends to $-\infty$.
Every hypothetical failure of the uniform upper limit
would have one of these two subsequences, so neither
is possible. Finally there are at most $2^{1+\delta n}$
conjugate terms, whose logarithm divided by $n^3$ tends
to zero. This proves the asserted uniform exponential-range
version of (P-F56as). The multipliers outside that range
remain an open part of (P-F56n).
\hfill$\square$
